\documentclass{report}
\usepackage{amsmath,amssymb}
\setlength{\parindent}{0mm}
\setlength{\parskip}{1em}
\begin{document}
\begin{center}
\rule{6in}{1pt} \
{\large Hormozd Gahvari \\
{\bf Performance Modeling in Algebraic Multigrid: Latest Developments}}

P O Box 808 \\ L-561 \\ Livermore \\ CA 94551
\\
{\tt gahvari1@llnl.gov}\\
William Gropp\\
Kirk E. Jordan\\
Martin Schulz\\
Ulrike Meier Yang\end{center}

Algebraic multigrid (AMG) solvers are very popular due to their ideal
algorithmic complexity and applicability to unstructured problems, but
there are concerns about scalability when solving large problems on
massively parallel machines. These concerns stem from large numbers of
messages being sent between processors on coarse grids. To improve the
performance and scalability of AMG in the parallel setting, we have
turned to performance modeling to guide data redistribution at runtime in
the AMG solve cycle that reduces the number of messages. A performance
model is employed on the fly to control when in the cycle the
redistribution occurs, and how much of it is performed. We are now in the
process of turning this inline modeling into a feature that can be
employed in production software. In this talk, we discuss what we have
achieved so far and our plans for the future. We have already been able
to obtain significant speedups on multiple machines on both simple and
more complicated problems, and we can now make machine parameter
measurements for the performance models on the fly instead of having to
input them beforehand.


\end{document}
